\documentclass[11pt,a4paper]{article}

\usepackage[T1]{fontenc}
\usepackage[utf8]{inputenc}
\usepackage{lmodern}
\usepackage{geometry}
\geometry{margin=2.5cm}

\usepackage{graphicx}
\usepackage{hyperref}
\usepackage{amsmath}
\usepackage{enumitem}
\usepackage{caption}

\usepackage{listings}
\usepackage{xcolor}

\definecolor{steelBlue}{rgb}{0.10, 0.25, 0.60}
\definecolor{slateGreen}{rgb}{0.20, 0.45, 0.30}
\definecolor{brickRed}{rgb}{0.70, 0.20, 0.20}
\definecolor{charcoal}{rgb}{0.15, 0.15, 0.15}
\definecolor{lightGrayBg}{rgb}{0.96, 0.96, 0.96}
\definecolor{frameGray}{rgb}{0.80, 0.80, 0.80}

\lstset{
    language=Java,
    backgroundcolor=\color{lightGrayBg},   
    commentstyle=\color{slateGreen}\itshape,   
    keywordstyle=\color{steelBlue}\bfseries,   
    stringstyle=\color{brickRed},              
    basicstyle=\ttfamily\small\color{charcoal},
    numberstyle=\tiny\color{gray},             
    breakatwhitespace=false,         
    breaklines=true,                 
    captionpos=b,                    
    keepspaces=true,                 
    numbers=left,                    
    numbersep=10pt,                  
    showspaces=false,                
    showstringspaces=false,
    showtabs=false,                  
    tabsize=4,
    frame=single,
    rulecolor=\color{frameGray}, 
    xleftmargin=25pt,
    framexleftmargin=25pt,
    framexrightmargin=5pt,
    aboveskip=1em,
    belowskip=1em
}

\hypersetup{
  colorlinks=true,
  linkcolor=blue!50!black,
  urlcolor=blue!50!black,
  citecolor=blue!50!black
}
\newcommand{\key}[1]{\fbox{\textsf{#1}}}

\newcommand{\course}{CmpE 160}
\newcommand{\assignment}{Assignment 1}
\newcommand{\studentName}{Name Surname}
\newcommand{\studentID}{2020XXXX}
\newcommand{\semester}{Spring 2026}
\newcommand{\sectionNo}{Section X}

\title{\course\  \assignment \\ \Large Report}
\author{\studentName\ \--- \studentID \\ \sectionNo \\ \semester}
\date{} % leave empty

\begin{document}
\maketitle

\vspace{-8pt}
\noindent\textbf{Note:} This report template is provided for convenience. You may modify the report structure as needed. Replace all placeholder text and figure names.

\section{Introduction}
Write a short introduction to the game and your implementation.
Briefly mention that the project is implemented in \textbf{Java} using the \textbf{StdDraw} library, and summarize the overall game.

\section{Game Dynamics}

\subsection{Stage 1: Main Menu}
Explain main menu dynamics and how you implemented. Include at least one screenshot.

\begin{figure}[ht!]
  \centering
  \includegraphics[width=0.75\linewidth]{example.png} % <-- replace filename
  \caption{Example caption.}
  \label{fig:example}
\end{figure}

\subsection{Stage 2: Gameplay}
Explain the gameplay such as player controls, shooting, enemies and how they move, scoring, drops (if any), hearts/lives, and level structure.
Include at least one gameplay screenshot.

\subsubsection{Enemies and Movement Mechanics}
Describe enemy formations and how enemies move.
If enemy movement uses randomness or different patterns across levels, describe the design briefly.

Below is an example of listlings package usage:

\begin{lstlisting}[caption={Example Function}, label={lst:enemy_move}]
private static double[] move(double[] pos, double[] vel) {
    // Update position
    double newPos = {pos[0]+vel[0], pos[1]+vel[1]};
    return newPos;
}
\end{lstlisting}
% --------------------------------

\subsubsection{Collision Detection}
Explain, at a high level, how collisions are detected (rectangle--rectangle overlap) and explain your code. You may include a debug screenshot showing bounding boxes.

\subsection{Stage 3: End-Game}
Explain the end-game scenarios:
(1) player wins, (2) player loses all lives (Game Over).
Explain the menu options on the end screen (Restart / End Game) and how selection works.
Include screenshots if available.

\end{document}